\documentclass[11pt,a4paper]{article}

\usepackage[margin=25mm]{geometry}
\usepackage[T1]{fontenc}
\usepackage{lmodern}
\usepackage{microtype}
\usepackage{amsmath,amssymb,bm}
\usepackage{booktabs,longtable,tabularx,array}
\usepackage{graphicx}
\usepackage{siunitx}
\usepackage[dvipsnames]{xcolor}
\usepackage{enumitem}
\usepackage{fancyhdr}
\usepackage{caption}
\usepackage{subcaption}
\usepackage[backend=biber,style=numeric-comp,sorting=nyt]{biblatex}
\usepackage{hyperref}
\usepackage[nameinlink,noabbrev]{cleveref}

\addbibresource{references.bib}
\graphicspath{{../figures/}{figures/}}
\hypersetup{colorlinks=true,linkcolor=MidnightBlue,citecolor=ForestGreen,urlcolor=BrickRed,pdftitle={SEMPHY01 Long Assessment 2}}
\sisetup{scientific-notation=true,separate-uncertainty=true,per-mode=symbol}
\setlength{\parindent}{0pt}
\setlength{\parskip}{0.6em}
\setlist{nosep}

\newcommand{\studentname}{TODO: Student Name}
\newcommand{\studentid}{TODO: Student ID}
\newcommand{\materialname}{TODO: Material and exact phase/polytype}
\newcommand{\submissiondate}{TODO: Submission date}
\newcommand{\repositoryurl}{https://example.invalid/replace-me}
\newcommand{\laonerevision}{TODO: LA1 commit, archive, or manifest hash}
\newcommand{\todo}[1]{\textcolor{BrickRed}{\textbf{TODO:} #1}}
\newcommand{\placeholderfigure}[3][0.86\linewidth]{\begin{figure}[htbp]\centering\fbox{\parbox[c][42mm][c]{#1}{\centering #2}}\caption{#2}\label{#3}\end{figure}}

\pagestyle{fancy}
\setlength{\headheight}{14pt}
\fancyhf{}
\lhead{SEMPHY01 --- Long Assessment 2}
\rhead{\studentname}
\cfoot{\thepage}

\begin{document}

\begin{titlepage}
\centering
{\Large De La Salle University\par}
\vspace{8mm}
{\large SEMPHY01 --- Semiconductor Physics\par}
\vspace{18mm}
{\Huge\bfseries Long Assessment 2\par}
\vspace{4mm}
{\LARGE From Equilibrium Carrier Statistics to Transport and p--n Junction Prediction\par}
\vspace{18mm}
\begin{tabular}{@{}rl@{}}
Student: & \studentname\\
Student ID: & \studentid\\
Material: & \materialname\\
Submission date: & \submissiondate\\
LA1 revision: & \laonerevision\\
Computational record: & \url{\repositoryurl}
\end{tabular}
\vfill
\begin{minipage}{0.88\textwidth}\small
This report must distinguish quantities inherited from LA1 from new mobility,
scattering, recombination, geometry, contact, and high-field assumptions.
Absolute mobility and lifetime are not inferred from effective mass and carrier
statistics alone.
\end{minipage}
\end{titlepage}

\tableofcontents
\clearpage

\section{Declaration, research question, and LA1 transfer}

\begin{table}[htbp]
\centering
\caption{Declared transport and device model.}
\begin{tabularx}{\textwidth}{@{}>{\bfseries}p{0.27\textwidth}X@{}}
\toprule
Item & Declaration\\
\midrule
Material continuity & \todo{Exact LA1 material, phase/polytype, and revision.}\\
Temperature range & \todo{Common validated range.}\\
Junction geometry & \todo{Orientation, area, doping, and neutral lengths.}\\
Mobility model & \todo{Electron/hole models, directions, and sources.}\\
Recombination model & \todo{Lifetime or microscopic model and injection range.}\\
Electrostatics & \todo{Depletion and numerical model.}\\
Current model & \todo{Finite/long base and nonideal terms.}\\
Dynamic prediction & \todo{Stored charge, switching, or collection calculation.}\\
Operating limit & \todo{Field, breakdown, heating, or documented exclusion.}\\
\bottomrule
\end{tabularx}
\end{table}

\subsection{Research question}
\todo{State one material-specific device prediction question. Compare notation
and junction conventions with established references such as Sze--Lee
\cite{sze-lee} and characterize measurements using an appropriate source such
as Schroder \cite{schroder}.}

\subsection{Transfer integrity}
\todo{Identify the immutable LA1 package, hashes, neutrality residuals, energy reference, and any corrected LA1 result.}

\placeholderfigure{LA1-to-LA2 model map distinguishing inherited state from new constitutive physics.}{fig:model-map}

\section{Homogeneous transport}

\subsection{Mobility, scattering, and provenance}
\todo{Define electron and hole mobility models. Explain why effective mass alone does not determine mobility.}

\begin{table}[htbp]
\centering
\caption{New transport and recombination parameters.}
\begin{tabularx}{\textwidth}{@{}llllX@{}}
\toprule
Quantity & Value/model & Unit & Valid range & Source and definition\\
\midrule
$\mu_n$ & \todo{ } & \si{\meter\squared\per\volt\per\second} & \todo{ } & \todo{ }\\
$\mu_p$ & \todo{ } & \si{\meter\squared\per\volt\per\second} & \todo{ } & \todo{ }\\
$\tau_n$ & \todo{ } & \si{\second} & \todo{ } & \todo{ }\\
$\tau_p$ & \todo{ } & \si{\second} & \todo{ } & \todo{ }\\
$\varepsilon_s$ & \todo{ } & \si{\farad\per\meter} & \todo{ } & \todo{ }\\
\bottomrule
\end{tabularx}
\end{table}

\subsection{Conductivity, resistivity, Hall response, and diffusion}
\begin{align}
\sigma&=q(n\mu_n+p\mu_p),\\
R_H&=\frac{p\mu_p^2-n\mu_n^2}{q(p\mu_p+n\mu_n)^2},\\
D_{n,p}&=\mu_{n,p}\frac{k_BT}{q}\quad\text{in the nondegenerate limit}.
\end{align}
\placeholderfigure{Mobility, conductivity/resistivity, Hall response, and Einstein-relation validity.}{fig:transport}
\todo{Trace each curve to the transferred LA1 carrier state and the new mobility model.}

\section{Recombination and continuity}

\todo{Define the lifetime or recombination model, injection range, and boundary physics.}
\begin{equation}
\frac{\partial\Delta c}{\partial t}=D\frac{\partial^2\Delta c}{\partial x^2}-v_d\frac{\partial\Delta c}{\partial x}-U(\Delta c)+G(x,t).
\end{equation}
\placeholderfigure{Analytic and finite-difference minority profiles, transient evolution, and conservation check.}{fig:continuity}

\begin{table}[htbp]
\centering
\caption{Continuity-solver validation.}
\begin{tabular}{@{}lrrr@{}}
\toprule
Check & Coarse & Refined & Interpretation\\
\midrule
Profile error & \todo{ } & \todo{ } & \todo{ }\\
Flux error & \todo{ } & \todo{ } & \todo{ }\\
Charge-balance residual & \todo{ } & \todo{ } & \todo{ }\\
Minimum concentration & \todo{ } & \todo{ } & \todo{ }\\
\bottomrule
\end{tabular}
\end{table}

\section{p--n-junction electrostatics and capacitance}

\subsection{Transferred bulk states and built-in potential}
\begin{equation}
qV_{\mathrm{bi}}=\mu_n^{\mathrm{bulk}}-\mu_p^{\mathrm{bulk}}.
\end{equation}
\todo{Compare chemical-potential, carrier-density, fully ionized, and numerical electrostatic results.}

\subsection{Depletion and numerical Poisson solutions}
\begin{equation}
W=\sqrt{\frac{2\varepsilon_s}{q}(V_{\mathrm{bi}}-V)\left(\frac{1}{N_A^-}+\frac{1}{N_D^+}\right)}.
\end{equation}
\placeholderfigure{Charge, field, potential, and band diagrams at equilibrium and selected biases.}{fig:electrostatics}
\placeholderfigure{Numerical Poisson result compared with the depletion approximation.}{fig:poisson}

\subsection{C--V prediction}
\begin{equation}
C_j=\frac{\varepsilon_s}{W}.
\end{equation}
\placeholderfigure{$C_j(V_R)$, $1/C_j^2$, and recovered ionized doping.}{fig:cv}

\section{Diode current, charge storage, and dynamics}

\subsection{Finite-base current}
\begin{equation}
J_s=q\left[\frac{D_nn_{p0}}{L_n}\coth\left(\frac{w_p}{L_n}\right)+\frac{D_pp_{n0}}{L_p}\coth\left(\frac{w_n}{L_p}\right)\right].
\end{equation}
\todo{Trace $n_{p0}$ and $p_{n0}$ to LA1 and $D$, $L$ to the new transport/recombination model.}
\placeholderfigure{Ideal and nonideal current components with extracted ideality factor.}{fig:iv}

\subsection{Stored charge and transient prediction}
\begin{align}
Q_s&=q\int\Delta c(x)\,\mathrm dx,\\
C_d&=\frac{\mathrm dQ_s}{\mathrm dV}.
\end{align}
\placeholderfigure{Stored charge, depletion and diffusion capacitance, and transient terminal response.}{fig:dynamics}

\subsection{Operating limit}
\todo{Estimate one field, breakdown, tunneling, velocity-saturation, punch-through, series-resistance, or heating limit.}

\section{Integrated prediction and uncertainty}

\begin{table}[htbp]
\centering
\caption{Declared device operating-point prediction.}
\begin{tabularx}{\textwidth}{@{}lrrX@{}}
\toprule
Quantity & Prediction & Uncertainty/range & Dominant source\\
\midrule
\todo{ } & \todo{ } & \todo{ } & \todo{ }\\
\todo{ } & \todo{ } & \todo{ } & \todo{ }\\
\bottomrule
\end{tabularx}
\end{table}

\placeholderfigure{Validity map in temperature, concentration, bias, field, injection, or length.}{fig:validity-map}

\section{Validation against independent evidence}
\todo{Separate calibration data from at least two independent validation comparisons.}

\section{Conclusion}
\todo{State the most reliable prediction, least reliable prediction, dominant LA1 parameter, dominant new LA2 uncertainty, first failing assumption, and next needed experiment or calculation.}

\printbibliography
\end{document}
